\documentclass[11pt,a4paper]{article}

% ------------------------------------------------------------------
%  Packages
% ------------------------------------------------------------------
\usepackage[T1]{fontenc}
\usepackage[utf8]{inputenc}
\usepackage{amsmath,amssymb,amsfonts}
\usepackage{mathtools}
\usepackage{booktabs}
\usepackage{array}
\usepackage{geometry}
\usepackage{xcolor}
\usepackage{enumitem}
\usepackage{hyperref}
\usepackage{graphicx}
\usepackage{caption}

\geometry{margin=2.5cm}
\emergencystretch=1.5em

\hypersetup{
  colorlinks=true,
  linkcolor=blue!60!black,
  citecolor=blue!60!black,
  urlcolor=blue!60!black,
  pdftitle={NS3D - Numerical methods and discretization},
  pdfauthor={J. Lux}
}

% ------------------------------------------------------------------
%  Notation helpers
% ------------------------------------------------------------------
\newcommand{\vc}[1]{\mathbf{#1}}            % vector (bold)
\newcommand{\tc}[1]{\mathbf{#1}}            % tensor (bold)
\newcommand{\RHS}{\mathrm{RHS}}
\newcommand{\dvc}[1]{\mathrm{d}\!#1}

% Face neighbours shorthand
\newcommand{\cw}{\mathrm{W}}
\newcommand{\ce}{\mathrm{E}}
\newcommand{\cn}{\mathrm{N}}
\newcommand{\cs}{\mathrm{S}}
\newcommand{\ct}{\mathrm{T}}
\newcommand{\cb}{\mathrm{B}}
\newcommand{\cc}{\mathrm{C}}

\title{\textbf{NS3D --- Numerical Methods and Discretization}\\[4pt]
\large Steady, low-Reynolds incompressible flow in 3-D porous media\\
\large (SIMPLE / SIMPLEC / SIMPLER / IDEAL / SIMPLER-R)}
\author{J.\ Lux}
\date{\today}

\begin{document}
\maketitle

\begin{abstract}
This document gives a self-contained description of the numerical methods
implemented in \texttt{NS3D}: the steady, low-Reynolds incompressible
Navier--Stokes equations discretized with a finite-volume method on a
staggered (MAC) grid, and the pressure--velocity coupling algorithms
\textsc{Simple}, \textsc{Simplec}, \textsc{Simpler}, \textsc{Ideal} and
\textsc{Simpler-R}.  The discrete momentum equation, the diffusive and
convective fluxes, the under-relaxation, the boundary conditions and the
derivation of the pressure correction are all written out explicitly.
The document is faithful to the reference C++ implementation in
\texttt{src/}.
\end{abstract}

\tableofcontents

% ==================================================================
\section{Governing equations}
% ==================================================================

The flow is assumed incompressible, steady and at low Reynolds number.
The continuity and momentum equations read
\begin{align}
\nabla\!\cdot\!\vc{u} &= 0 , \label{eq:cont} \\[4pt]
\rho\,\vc{u}\!\cdot\!\nabla\,\vc{u}
&= -\nabla p + \mu\,\nabla^{2}\vc{u} , \label{eq:mom}
\end{align}
where $\vc{u}=(u,v,w)$ is the velocity, $p$ the pressure, $\rho$ the
density and $\mu$ the dynamic viscosity.  In component form the $x$-momentum
equation is
\begin{equation}
\mu\left(\frac{\partial^{2}u}{\partial x^{2}}
      +\frac{\partial^{2}u}{\partial y^{2}}
      +\frac{\partial^{2}u}{\partial z^{2}}\right)
 -\rho\left(u\frac{\partial u}{\partial x}
          +v\frac{\partial u}{\partial y}
          +w\frac{\partial u}{\partial z}\right)
 = \frac{\partial p}{\partial x} ,
\end{equation}
with analogous expressions for the $y$ and $z$ components.

% ==================================================================
\section{Discretization}
% ==================================================================

\subsection{Staggered (MAC) grid}

A staggered grid is used: the pressure $p$ is stored at cell centres and the
three velocity components are stored on the cell faces, one component per
face.  A pressure node $C$ has six face neighbours, denoted
$\cw,\ce$ (along $x$), $\cs,\cn$ (along $y$) and $\cb,\ct$ (along $z$).
On the face between $C$ and $\ce$ the $u$-component lives, on the face between
$C$ and $\cn$ the $v$-component lives, and so on.  The mesh is assumed
Cartesian and locally uniform with spacings $\Delta x,\Delta y,\Delta z$.

\subsection{Finite-volume form}

Equation \eqref{eq:mom} is integrated over a control volume $V$ bounded by
$S$.  Applying the divergence theorem to the convection and diffusion terms
gives, for the $u$-component,
\begin{equation}
\frac{1}{V}\!\int_{S}\!\mu\,\nabla u\!\cdot\!\vc{n}\,\dvc{S}
-\frac{1}{V}\!\int_{S}\!\rho\,(\vc{u}\!\cdot\!\vc{n})\,u\,\dvc{S}
=\frac{1}{V}\!\int_{V}\frac{\partial p}{\partial x}\,\dvc{V} .
\end{equation}
Writing the surface integrals as the sum of the contributions of the six
faces, the discrete equation for the $u$-node $C$ takes the standard linear
form
\begin{equation}
a^{\,u}_{\cc}\,u_{\cc}+\sum_{n\in\{\cw,\ce,\cs,\cn,\cb,\ct\}} a^{\,u}_{n}\,u_{n}
=\frac{p_{\cc}-p_{\cw}}{\Delta x}+b ,
\label{eq:discrete-u}
\end{equation}
where the coefficients $a^{\,u}_{n}$ are built from the diffusive and
convective fluxes on face $n$ (Sections~\ref{sec:diffusive}--\ref{sec:convective}),
and the right-hand
side is the imposed pressure gradient (or the boundary value).  Equations of
the same form hold for $v$ (pressure gradient $p_{\cc}-p_{\cs}$ over
$\Delta y$) and $w$ (over $\Delta z$).

\subsection{Diffusive fluxes}
\label{sec:diffusive}

The diffusive contribution on a face is a second-order central difference.
For the $u$-node $C$:
\begin{align}
D^{\,u}_{\cw}&=\mu\frac{u_{\cc}-u_{\cw}}{\Delta x^{2}},
&
D^{\,u}_{\ce}&=\mu\frac{u_{\ce}-u_{\cc}}{\Delta x^{2}},\\[2pt]
D^{\,u}_{\cs}&=\mu\frac{u_{\cc}-u_{\cs}}{\Delta y^{2}},
&
D^{\,u}_{\cn}&=\mu\frac{u_{\cn}-u_{\cc}}{\Delta y^{2}},\\[2pt]
D^{\,u}_{\cb}&=\mu\frac{u_{\cc}-u_{\cb}}{\Delta z^{2}},
&
D^{\,u}_{\ct}&=\mu\frac{u_{\ct}-u_{\cc}}{\Delta z^{2}} .
\end{align}

\subsection{Convective fluxes}
\label{sec:convective}

Let $u^{*},v^{*},w^{*}$ denote the velocity field from the previous iteration
(the coefficients are lagged).  The convective flux on a face is either
\emph{centered} or \emph{upwind}, selected by a local P\'eclet criterion.
Writing $D=\mu/\Delta x^{2}$ (or $\mu/\Delta y^{2}$, $\mu/\Delta z^{2}$) and
$F$ for the signed face flux, the switch implemented in
\texttt{NSBase::AssembleMomentum\_*} is
\begin{equation}
|F|\,\Delta x^{2} \;>\; 2\,\mu\,s_{\mathrm{sch}}
\quad\Longrightarrow\quad \text{upwind},
\qquad\text{otherwise centered},
\end{equation}
where $s_{\mathrm{sch}}$ is the \emph{scheme factor}.  With the default
$s_{\mathrm{sch}}=0$ the test reduces to $|F|>0$, i.e.\ the upwind scheme is
always used whenever the face flux is non-zero; the centered scheme is taken
only for a strictly zero flux.  This is the intended behaviour for the
low-Reynolds flows solved here (an always-upwind convection is
unconditionally stable).

The upwind coefficients on a face $n$ are
\begin{equation}
a^{\,u}_{\cc}\;\mathrel{+}=\;-D+\min(F,0),
\qquad
a^{\,u}_{n}\;=\;D+\max(F,0),
\label{eq:upwind}
\end{equation}
and the centered coefficients are
\begin{equation}
a^{\,u}_{\cc}\;\mathrel{+}=\;-D+\tfrac{1}{2}F,
\qquad
a^{\,u}_{n}\;=\;D+\tfrac{1}{2}F .
\label{eq:centered}
\end{equation}
Here $F$ is the signed convective flux (positive for inflow through the face).
For the two faces along the $u$-axis the flux is built from the $u$-values on
either side of the face; for the transverse faces the face velocity is the
average of the two $v$ (resp.\ $w$) values flanking the face, e.g.\ for the
$\cn$ face of node $C$, $v_{\cn}=\tfrac12\bigl(v_{\cc-}+v_{\cc}\bigr)$ and
$F=-v_{\cn}\,\rho/\Delta y$.

\subsection{Assembled coefficients}
\label{sec:coeffs}

For the $u$-node $C$ the diagonal coefficient is the sum over the six faces
of the central part,
\begin{equation}
a^{\,u}_{\cc}=\sum_{n} D^{\,u}_{n}\;-\;\sum_{n}\min(F_{n},0)
\qquad(\text{upwind}) ,
\end{equation}
and each off-diagonal is given by \eqref{eq:upwind} or \eqref{eq:centered}.
The $v$ and $w$ systems are identical with the indices shifted.  The three
momentum systems are assembled as sparse matrices $\mathbf{A}^{u},
\mathbf{A}^{v},\mathbf{A}^{w}$ and factorized with an iterative linear
solver.

% ==================================================================
\section{Under-relaxation}
% ==================================================================

The $u$-update is relaxed with a factor $\alpha\in(0,1]$.  Writing the
unrelaxed equation \eqref{eq:discrete-u} as
\begin{equation}
a^{\,u}_{\cc}\,u_{\cc}^{\,n}+\sum_{n}a^{\,u}_{n}\,u_{n}^{\,n}
=\frac{p_{\cc}-p_{\cw}}{\Delta x}+b ,
\end{equation}
and imposing $u^{\,n}-u^{\,n-1}=\alpha\,(\text{correction})$, one obtains
the relaxed system solved by the code:
\begin{equation}
\boxed{\;
\frac{a^{\,u}_{\cc}}{\alpha}\,u_{\cc}^{\,n}+\sum_{n}a^{\,u}_{n}\,u_{n}^{\,n}
=\frac{p_{\cc}-p_{\cw}}{\Delta x}+b
+\frac{1-\alpha}{\alpha}\,a^{\,u}_{\cc}\,u_{\cc}^{\,n-1}
\;}
\label{eq:relaxed}
\end{equation}
In the implementation only the diagonal is scaled by $1/\alpha$
(\texttt{A(c,c) /= alpha}); the right-hand side is augmented with
\texttt{u(c) * A(c,c) * (1-alpha)} (see
\texttt{Assemble\_RHS\_*}).  A separate pressure under-relaxation factor
$\beta$ is applied to the pressure equation and to the pressure update
(Section~\ref{sec:coupling}).

% ==================================================================
\section{Boundary conditions}
% ==================================================================

\begin{itemize}[leftmargin=1.5em]
\item \textbf{No-slip solid/fluid interface.}  The velocity is zero at the
  face.  The convective flux vanishes, but the diffusive flux is retained with
  a halved distance, which doubles the diagonal contribution on that face:
  \begin{equation}
  a^{\,u}_{\cc}\;\mathrel{+}=\;-2D
  \qquad(\text{solid face}).
  \end{equation}

\item \textbf{Inlet.}  The velocity is imposed; only the component along the
  flow direction is non-zero.  The pressure is fixed to $p_{0}$, so the
  right-hand side reads $(p_{\cc}-p_{0})/\Delta x$.

\item \textbf{Outlet.}  The velocity is extrapolated (zero gradient) and the
  pressure is fixed to $p_{1}$; the right-hand side reads
  $(p_{1}-p_{\cw})/\Delta x$.
\end{itemize}

% ==================================================================
\section{Pressure--velocity coupling}
% ==================================================================
\label{sec:coupling}

\subsection{Derivation of the correction relation}

The momentum equation written for a \emph{guess} velocity $\vc{u}^{*}$ is
\begin{equation}
\frac{a^{\,u}_{\cc}}{\alpha}\,u^{*}_{\cc}
+\sum_{n}a^{\,u}_{n}\,u^{*}_{n}
=\frac{p^{(n-1)}_{\cc}-p^{(n-1)}_{\cw}}{\Delta x}
+b+\frac{1-\alpha}{\alpha}\,a^{\,u}_{\cc}\,u^{(n-1)}_{\cc} .
\label{eq:guess}
\end{equation}
Define the \emph{pseudo-velocity} (or \emph{H-function})
\begin{equation}
\widehat{u}_{\cc}
=-\frac{\sum_{n}a^{\,u}_{n}\,u^{*}_{n}}{a^{\,u}_{\cc}/\alpha}
+\frac{1-\alpha}{\alpha}\,a^{\,u}_{\cc}\,u^{(n-1)}_{\cc}
\;=\;
-\frac{\alpha}{a^{\,u}_{\cc}}\sum_{n}a^{\,u}_{n}\,u^{*}_{n}
+(1-\alpha)\,u^{(n-1)}_{\cc} ,
\label{eq:uhat}
\end{equation}
so that \eqref{eq:guess} becomes the simple, pressure-only relation
\begin{equation}
\boxed{\;
u^{*}_{\cc}=\widehat{u}_{\cc}
+\frac{1}{a^{\,u}_{\cc}/\alpha}\,
\frac{p_{\cc}-p_{\cw}}{\Delta x}
\;}
\label{eq:urel}
\end{equation}
which the code computes in \texttt{Compute\_\{u,v,w\}\_hat} and
\texttt{CorrectVelocityField\allowbreak SIMPLER}.

The actual (corrected) velocity $\vc{u}=\vc{u}^{*}+\vc{u}^{\prime}$ is
required to satisfy both the momentum equation and the continuity equation.
Subtracting \eqref{eq:guess} from the equation for $\vc{u}$ and using
$\nabla\!\cdot\!\vc{u}^{\prime}=\nabla\!\cdot\!\vc{u}=0$ gives, to first
order,
\begin{equation}
\frac{a^{\,u}_{\cc}}{\alpha}\,u^{\prime}_{\cc}
+\sum_{n}a^{\,u}_{n}\,u^{\prime}_{n}
=\frac{p^{\prime}_{\cc}-p^{\prime}_{\cw}}{\Delta x} .
\label{eq:pcorr}
\end{equation}
Each algorithm differs in how the sum over the neighbours
$\sum_{n}a^{\,u}_{n}u^{\prime}_{n}$ in \eqref{eq:pcorr} is approximated.

\subsection{\textsc{Simple}}
\label{sec:simple}

The \textsc{Simple} algorithm (Semi-Implicit Method for Pressure Linked
Equations) was introduced by Patankar and Spalding
\cite{patankar1972,patankar1980}.  The neighbour correction is neglected,
$\sum_{n}a^{\,u}_{n}u^{\prime}_{n}\approx 0$, so that the face correction is
\begin{equation}
u^{\prime}_{\cc}=\frac{1}{a^{\,u}_{\cc}/\alpha}\,
\frac{p^{\prime}_{\cc}-p^{\prime}_{\cw}}{\Delta x}
\qquad\Longleftrightarrow\qquad
u^{\prime}_{\cc}=\frac{\alpha}{a^{\,u}_{\cc}}\,
\frac{p^{\prime}_{\cc}-p^{\prime}_{\cw}}{\Delta x} .
\label{eq:simple}
\end{equation}
Substituting \eqref{eq:simple} into the continuity equation
$\nabla\!\cdot\!\vc{u}^{*}=0$ yields the pressure-correction Poisson system.
With
\begin{equation}
d^{\,u}_{\cc}=\frac{1}{(a^{\,u}_{\cc}/\alpha)\,\Delta x^{2}} ,
\end{equation}
the node equation is
\begin{equation}
\sum_{n}\Bigl[d^{\,u}_{\cc}p^{\prime}_{\cc}-d^{\,u}_{\cc}p^{\prime}_{n}\Bigr]
=\sum_{n}\frac{u^{*}_{\text{face}}}{\Delta x} ,
\end{equation}
i.e.\ a five-point (six-point in 3-D) Laplacian with coefficients
$d^{\,u}_{\cc},d^{\,v}_{\cc},d^{\,w}_{\cc}$ and a source made of the
divergence of the guessed velocity.  This is
\texttt{AssemblePressureCorrectionSIMPLE} and
\texttt{CorrectVelocityField\allowbreak SIMPLE}.

\subsection{\textsc{Simplec}}
\label{sec:simplec}

\textsc{Simplec} (SIMPLE-consistent) was proposed by Manasria and Mudawar
\cite{manasria1995}.  The neglected sum in \eqref{eq:pcorr} is approximated
by the central value,
\begin{equation}
\sum_{n}a^{\,u}_{n}u^{\prime}_{n}\;\approx\;
u^{\prime}_{\cc}\sum_{n}a^{\,u}_{n} ,
\end{equation}
which gives
\begin{equation}
\boxed{\;
u^{\prime}_{\cc}=
\frac{1}{\dfrac{a^{\,u}_{\cc}}{\alpha}+\displaystyle\sum_{n}a^{\,u}_{n}}\,
\frac{p^{\prime}_{\cc}-p^{\prime}_{\cw}}{\Delta x}
\;}
\label{eq:simplec}
\end{equation}
Note that
$\dfrac{a^{\,u}_{\cc}}{\alpha}+\sum_{n}a^{\,u}_{n}$ is exactly the
\emph{row sum} of the momentum matrix row $\cc$, which is what the code uses
(\texttt{A.row(c).sum()}).  The pressure-correction coefficients are therefore
\begin{equation}
d^{\,u}_{\cc}
=\frac{1}{\Delta x^{2}}\,\left[\,\frac{a^{\,u}_{\cc}}{\alpha}+\sum_{n}a^{\,u}_{n}\,\right]^{-1} ,
\end{equation}
assembled in \texttt{AssemblePressure\allowbreak Correction\allowbreak SIMPLEC}
and applied in
\texttt{CorrectVelocityField\allowbreak SIMPLEC}.

\subsection{\textsc{Simpler}}
\label{sec:simpler}

\textsc{Simpler} (SIMPLE Reformulation) was introduced by Acharya
\cite{acharya1991}.  It splits the update into a pressure step and a velocity
correction step, using the pseudo-velocity \eqref{eq:urel}.  One outer
iteration is:
\begin{enumerate}[leftmargin=1.5em]
\item \textbf{Pressure step} (inner loop, $N_{1}$ times): from the current
  velocity compute the pseudo-velocities $\widehat{u},\widehat{v},\widehat{w}$
   (\eqref{eq:uhat}) and solve the continuity system (with the
   \textsc{simple} coefficients of \eqref{eq:simple}) for a new pressure
   $p^{*}$.  Then update
   the working velocity using \eqref{eq:urel}
   (\texttt{AssemblePressure} + \texttt{CorrectVelocityField\allowbreak SIMPLER}).
\item \textbf{Velocity step}: solve the relaxed momentum system
  \eqref{eq:relaxed} with the new pressure $p^{*}$ to obtain the velocity
  $\vc{u}^{*}$ (\texttt{ComputeVelocityField}).
\item \textbf{Correction step}: assemble and solve the \textsc{simple}
  pressure-correction system \eqref{eq:simple} for $p^{\prime}$ and apply
  \eqref{eq:simple} to the velocity only
  (\texttt{CorrectVelocityField\allowbreak SIMPLE}).  The pressure itself is
  \emph{not} updated with $p^{\prime}$.
\end{enumerate}
Extra memory: the pseudo-velocities $\widehat{u},\widehat{v},\widehat{w}$.

\subsection{\textsc{Ideal}}
\label{sec:ideal}

\textsc{Ideal} (Inner Doubly Iterative Efficient Algorithm for Linked
Equations) was proposed by Sun \emph{et al.} \cite{sun2007}.  It keeps two
pressure solves and one momentum solve, and updates
the pressure at the end of the iteration (unlike \textsc{Simpler}).  One outer
iteration is:
\begin{enumerate}[leftmargin=1.5em]
\item \textbf{Pressure step 1} (inner loop, $N_{1}$ times): compute
  pseudo-velocities from the working velocity, solve the continuity system for
  $p^{*}$ and update the working velocity with \eqref{eq:urel}.
\item \textbf{Velocity step}: solve the relaxed momentum system with $p^{*}$
  to obtain $\vc{u}^{*}$.
\item \textbf{Pressure step 2} (inner loop, $N_{2}$ times): compute the
  pseudo-velocities from $\vc{u}^{*}$, solve the continuity system again (this
  time \emph{without} pressure under-relaxation) for $p^{**}$ and update the
  velocity with \eqref{eq:urel}.  The final pressure is $p^{**}$.
\end{enumerate}
The advantage over \textsc{Simpler} is that the final pressure is obtained by
a direct continuity solve rather than by a pressure-correction solve, which is
typically better conditioned.

\subsection{\textsc{Simpler-R}}
\label{sec:simplerr}

\textsc{Simpler-R} (SIMPLER-revised) \cite{li2017} is a hybrid of
\textsc{Ideal} and \textsc{Simpler}:
\begin{enumerate}[leftmargin=1.5em]
\item \textbf{Pressure step 1} (inner loop, $N_{1}$ times): as in
  \textsc{Ideal}, solve the continuity system for $p^{*}$ and update the
  working velocity with \eqref{eq:urel}.
\item \textbf{Velocity step}: solve the relaxed momentum system with $p^{*}$
  to obtain $\vc{u}^{*}$.
\item \textbf{Correction step}: solve the \textsc{simple} pressure-correction
  system \eqref{eq:simple} for $p^{\prime}$ and apply \eqref{eq:simple} to the
  velocity.
\item \textbf{Pressure refresh}: recompute the pseudo-velocities from the
  corrected velocity, then solve the continuity system once more (with pressure
  under-relaxation) to update the pressure $p$.
\end{enumerate}
This combines the velocity correction of \textsc{Simpler} with the pressure
refresh of \textsc{Ideal}.

\subsection{Summary of the schemes}
\label{sec:summary}

A direct numerical comparison of \textsc{Piso}, \textsc{Simpler} and
\textsc{Simplec} for steady incompressible flow was performed by Jang \emph{et
al.} \cite{jang1986}; the relative cost and conditioning of the pressure solves
summarized below is consistent with that study.

\begin{table}[h]
\centering
\small
\begin{tabular}{@{}lllll@{}}
\toprule
Scheme & \# pressure solves & Momentum solve & Velocity correction & Pressure corrected\\
\midrule
\textsc{Simple}    & 1 & 1 & \eqref{eq:simple}  (diag.)  & yes ($\beta\,p^{\prime}$)\\
\textsc{Simplec}   & 1 & 1 & \eqref{eq:simplec} (row sum) & yes ($\beta\,p^{\prime}$)\\
\textsc{Simpler}   & $N_{1}+1$ & 1 & \eqref{eq:simple} (diag.) & no\\
\textsc{Ideal}     & $N_{1}+N_{2}$ & 1 & \eqref{eq:urel} & yes (direct)\\
\textsc{Simpler-R} & $N_{1}+2$ & 1 & \eqref{eq:simple} (diag.) & yes (direct)\\
\bottomrule
\end{tabular}
\caption{Comparison of the five pressure--velocity coupling schemes.  The
``\# pressure solves'' counts the continuity solves per outer iteration
(including the inner $N_{1}$/$N_{2}$ loops).}
\label{tab:schemes}
\end{table}

% ==================================================================
\section{Convergence control}
% ==================================================================

Two mechanisms, both implemented in
\texttt{NavierStokesBase::RelaxControl} and the velocity-correction
functions, are used to keep the iteration stable.

\subsection{Adaptive relaxation}

Let $\mathrm{RMR}$ be the relative mass residual.  The relaxation factors are
adapted as follows:
\begin{itemize}[leftmargin=1.5em]
\item \textbf{Divergence} ($\mathrm{RMR}>\mathrm{RMR}_{\text{old}}\,\delta_{\max}$):
  decrease $\alpha$ (and $\beta$ unless it is held constant),
  $\alpha\leftarrow\max(\alpha_{\min},\alpha-\Delta)$, and reset
  $\mathrm{RMR}_{\text{old}}\leftarrow 1.1\,\mathrm{RMR}_{\text{old}}$.
\item \textbf{Oscillation} (the recent average residual is flat): decrease
  $\alpha,\beta$ by a smaller increment ($\Delta/4$).
\item \textbf{Slow divergence}: decrease $\alpha,\beta$ by $\Delta$.
\end{itemize}

\subsection{Velocity limiting}

After each velocity solve/correction, the per-node velocity change is clamped
to a maximum increment $\delta v_{\max}$.  In the code
$\delta v_{\max}=v_{\text{moy}}\min(1,\alpha\,1.5)$ for the momentum solve and
a large value ($100\,v_{\text{moy}}$, i.e.\ effectively unbounded) for the
correction steps, so that only the main momentum update is limited.

% ==================================================================
\section{Derived quantities}
% ==================================================================

\subsection{Permeability (Darcy)}

For a fully-developed, unidirectional flow the Darcy law
$-\nabla\langle p\rangle=\mu\,\mathbf{K}^{-1}\langle\vc{u}\rangle$ is used to
define the permeability from the computed mean velocity $\langle v\rangle$ and
the imposed pressure difference:
\begin{equation}
K=\frac{\langle v\rangle\,\mu\,\Delta x\,\bigl(N_{x}-2\,b\bigr)}
        {p_{\text{in}}-p_{\text{out}}} ,
\end{equation}
where $N_{x}$ is the number of cells along the flow direction and $b$ the
number of solid border cells.  The mean velocity is obtained by averaging the
velocity over the inlet/outlet faces.  The code guards against division by
zero when $p_{\text{in}}=p_{\text{out}}$.

\subsection{Forchheimer}

For flows where inertia is non-negligible, the macroscopic momentum balance is
the Forchheimer (or Darcy--Forchheimer) relation
\begin{equation}
-\nabla\langle p\rangle
=\frac{\mu}{k_{1}}\,\langle v\rangle
 +\frac{\rho}{k_{2}}\,\langle v\rangle^{2} ,
\end{equation}
with the linear coefficient $k_{1}$ (permeability) and the inertial
coefficient $k_{2}$.  The two coefficients are identified by fitting the
computed $(\Delta p,\langle v\rangle)$ data to this quadratic.

\subsection{Dimensionless permeability}

Defining the reference quantities $L_{\text{ref}}$ (length),
$v_{\text{ref}}$ (velocity) and
$\Delta\langle P^{*}\rangle$ (dimensionless pressure drop), the dimensional
permeability is
\begin{equation}
K=\frac{\langle v^{*}\rangle\,\Delta L^{*}}{\epsilon\,\Delta\langle P^{*}\rangle}
   \;\frac{\mu\,L_{\text{ref}}}{\rho\,v_{\text{ref}}}
=\frac{L_{\text{ref}}^{2}}{\mathrm{Re}}\;K^{*} ,
\end{equation}
where $\mathrm{Re}=\rho v_{\text{ref}}L_{\text{ref}}/\mu$ and $\epsilon$ the
porosity.

% ==================================================================
\section*{Notation}
% ==================================================================
\addcontentsline{toc}{section}{Notation}
\begin{center}
\small
\begin{tabular}{@{}ll@{}}
\toprule
Symbol & Meaning\\
\midrule
$\vc{u}=(u,v,w)$ & velocity vector\\
$p$ & pressure\\
$\rho,\mu$ & density, dynamic viscosity\\
$\alpha,\beta$ & velocity, pressure under-relaxation factors\\
$\cc,\cw,\ce,\cs,\cn,\cb,\ct$ & cell centre and face neighbours\\
$a^{\,u}_{\cc}$ & diagonal coefficient of the $u$-momentum equation\\
$a^{\,u}_{n}$ & off-diagonal coefficient towards neighbour $n$\\
$D^{\,u}_{n}$ & diffusive flux on face $n$\\
$F_{n}$ & signed convective flux on face $n$\\
$\widehat{u}$ & pseudo-velocity (H-function)\\
$p^{\prime}$ & pressure correction\\
$d^{\,u}_{\cc}$ & pressure-correction coefficient for $u$\\
$N_{1},N_{2}$ & inner-loop counts (\textsc{Simpler}/\textsc{Ideal})\\
$\mathrm{RMR}$ & relative mass residual\\
$K,k_{1},k_{2}$ & permeability, Forchheimer coefficients\\
$\Delta x,\Delta y,\Delta z$ & mesh spacings\\
\bottomrule
\end{tabular}
\end{center}

% ==================================================================
%  thebibliography already prints a `\section*{\refname}' heading, so no
%  explicit heading is needed here.
\addcontentsline{toc}{section}{References}
\begin{thebibliography}{9}

\bibitem{patankar1972}
S.\,V.\ Patankar and D.\,B.\ Spalding, ``A calculation procedure for heat,
mass and momentum transfer in three-dimensional parabolic flows,''
\emph{Int.\ J.\ Heat Mass Transfer}, vol.~15, no.~10, pp.~1787--1806, 1972.
doi:~\url{10.1016/0017-9310(72)90054-3}.

\bibitem{patankar1980}
S.\,V.\ Patankar, \emph{Numerical Heat Transfer and Fluid Flow}, Hemisphere
Publishing, Washington, D.C., 1980.

\bibitem{manasria1995}
R.\ Manasria and I.\ Mudawar, ``Analysis of the accuracy of the SIMPLE,
SIMPLER, and SIMPLEC algorithms for the calculation of natural and forced
convection,'' \emph{ASME J.\ Heat Transfer}, vol.~117, no.~1, pp.~157--163,
1995.

\bibitem{acharya1991}
C.\,A.\ Acharya, ``SIMPLER: A Simple Pressure Reformulation for the
Calculation of Incompressible Flows,'' \emph{ASME J.\ Fluids Eng.}, vol.~113,
pp.~292--296, 1991.

\bibitem{sun2007}
D.\,L.\ Sun, Z.\,G.\ Qu, Y.\,L.\ He, and W.\,Q.\ Tao, ``An Efficient
Segregated Algorithm for Incompressible Fluid Flow and Heat Transfer
Problems---IDEAL (Inner Doubly Iterative Efficient Algorithm for Linked
Equations) Part I: Mathematical Formulation and Solution Procedure,''
\emph{Numer.\ Heat Transfer, Part B: Fundam.}, vol.~53, no.~1, pp.~1--17,
2007. doi:~\url{10.1080/10407790701632543}.

\bibitem{li2017}
J.\ Li, Q.\ Zhang, and Z.\,-Q.\ Zhai, ``An efficient SIMPLER-revised
algorithm for incompressible flow with unstructured grids,''
\emph{Numer.\ Heat Transfer, Part B: Fundam.}, vol.~71, no.~5, pp.~425--442,
2017. doi:~\url{10.1080/10407790.2017.1293965}.

\bibitem{jang1986}
D.\,S.\ Jang, R.\ Jetli, and S.\ Acharya, ``Comparison of the PISO, SIMPLER,
and SIMPLEC algorithms for the treatment of the pressure--velocity coupling
in steady flow problems,'' \emph{Numer.\ Heat Transfer, Part B: Fundam.},
vol.~10, no.~3, pp.~209--228, 1986. doi:~\url{10.1080/10407798608552506}.

\end{thebibliography}

\end{document}
